\section{Introduction}

Spatially explicit intra-urban population surfaces are useful for service planning, hotspot screening, exposure assessment, urban comparison, and first-pass neighborhood analysis. Many practical decisions depend not only on total city population but also on how population is distributed across the urban fabric. Yet in data-scarce settings, the population is usually observed only through aggregated administrative counts rather than through true cell-level census labels \cite{Leyk2019,Mennis2009,Stevens2015}. This creates an immediate methodological asymmetry: the spatial product that urban analysis needs is finer than the population data that are actually observed.

In Kazakhstan and similar contexts, this asymmetry is especially important. Official city totals may exist and can be verified, but fine-resolution cell-level truth is unavailable for direct supervised learning. At the same time, open geospatial data such as buildings, roads, and points of interest provide a reproducible but indirect description of the urban environment. These signals do not reveal the true population count in each cell, but they can support a bounded proxy modeling workflow when linked to official totals and explicit allocation logic.

This problem has two linked parts. First, a reproducible calibrated proxy population surface must be constructed and validated under missing cell-level truth. Second, once that surface exists, users still need to know where it appears more stable and where it should be treated more cautiously. A deterministic surface that assigns one value per cell can otherwise encourage over-reading by making all cells look equally trustworthy. The scientific need is therefore not only for a calibrated proxy surface, but for a reliability-aware interpretation layer around that surface. Put differently, Stage 1 answers whether a bounded proxy surface can be built at all, while Stage 2 answers how that surface should be read once it exists.

City1 addresses these two research questions as one framework paper:

\begin{itemize}
\item \textbf{RQ1:} How can a reproducible calibrated proxy population surface be constructed under missing cell-level census truth?
\item \textbf{RQ2:} Once such a proxy surface is constructed, how can uncertainty, confidence, and hotspot stability be represented so that users do not over-read deterministic cell-level patterns?
\end{itemize}

The framework follows one mandatory transition logic. First, it constructs and validates a deterministic calibrated proxy surface under missing cell-level truth. Second, it extends this surface with uncertainty-aware interpretation, so that the output is not only spatially explicit but also reliability-aware.

\paragraph{Unified contributions.}
\begin{itemize}
\item a reproducible open-data proxy population surface pipeline based on weak supervision, Random Forest modeling, a fixed 500 m grid, and official-total calibration;
\item a multi-level validation stack for structural credibility under missing cell-level truth;
\item an uncertainty-aware extension built from calibrated ensemble intervals, relative uncertainty, and bounded confidence bands;
\item confidence-aware hotspot prioritization that separates stronger screening candidates from review-oriented ones;
\item an integrated evidence package in which figures, tables, formulas, limitations, and supplement material remain traceable to reference artifacts rather than to moving experimental targets.
\end{itemize}

Table~\ref{tab:contribution-map} summarizes these contributions as one two-stage framework rather than as two pasted papers. It is the roadmap for the paper, not a decorative summary.

\begin{table}[H]
\centering
\scriptsize
\caption{Contribution Map for the unified City1 framework. The table makes explicit how the calibrated baseline stage and the uncertainty-aware stage connect within one paper.}
\label{tab:contribution-map}
\begin{tabular}{p{3.0cm}p{1.2cm}p{3.0cm}p{3.0cm}p{3.8cm}}
\toprule
Contribution & Stage & Method / artifact & Main evidence & Claim boundary \\
\midrule
Reproducible calibrated proxy population surface & Stage 1 & weak target, Random Forest, 500 m grid, official-total calibration & LOCO, spatial block CV, grid benchmark, external benchmark, ablation & supports proxy-surface credibility, not cell-level census truth \\
Multi-level structural validation & Stage 1 & district benchmark, WorldPop/GHS-POP comparison, qualitative plausibility & district, external, and qualitative evidence & supports structural plausibility, not ground-truth validation \\
Uncertainty-aware interpretation layer & Stage 2 & calibrated ensemble, P10/P50/P90, interval width, relative uncertainty & interval coverage and error-width alignment & supports bounded uncertainty behavior, not true census uncertainty \\
Confidence-aware hotspot prioritization & Stage 2 & confidence bands and hotspot classes & confidence-band separation and hotspot stability & supports screening reliability, not verified population-hotspot truth \\
Tool-grounded interpretation layer & Stage 3 & read-only evidence tools, deterministic fallback, optional provider generation, cache support, and guardrails & bounded evaluation of evidence-linked explanation and safety behavior & supports reviewer-safe interpretation, not new population estimation or automated policy evidence \\
\bottomrule
\end{tabular}
\end{table}


The scientific identity of the paper is intentionally bounded. City1 does not reconstruct true cell-level census counts. It does not estimate true census uncertainty. Instead, it produces an officially calibrated proxy population surface and then adds a reliability-aware interpretation layer so that the resulting surface can be used more responsibly for screening and exploratory urban analysis. Within those bounds, the contribution is not merely a map. The contribution is a disciplined framework for producing, validating, and interpreting a calibrated proxy population surface under missing cell-level truth.

The rest of the manuscript follows this roadmap in order: Stage 1 constructs and validates the calibrated proxy baseline, Stage 2 explains how to interpret that baseline cautiously, and the discussion then connects the two stages into one coherent framework.
